\documentclass[11pt, a4paper]{article}

\usepackage[margin=2.5cm]{geometry}
\usepackage[T1]{fontenc}
\usepackage{lmodern}
\usepackage{microtype}
\usepackage{parskip}
\usepackage{titlesec}
\usepackage{xcolor}
\usepackage{listings}
\usepackage{mdframed}
\usepackage{booktabs}
\usepackage{hyperref}
\usepackage{enumitem}
\usepackage{fancyhdr}
\usepackage{tcolorbox}

% ---- colours ----------------------------------------------------------------
\definecolor{codebg}{HTML}{F5F5F5}
\definecolor{codeframe}{HTML}{CCCCCC}
\definecolor{keywordcol}{HTML}{0060AD}
\definecolor{commentcol}{HTML}{5C8A00}
\definecolor{stringcol}{HTML}{C0392B}
\definecolor{numbercol}{HTML}{7D3C98}
\definecolor{notebg}{HTML}{EAF4FB}
\definecolor{noteframe}{HTML}{5DADE2}
\definecolor{warnbg}{HTML}{FEF9E7}
\definecolor{warnframe}{HTML}{F39C12}

% ---- code listings ----------------------------------------------------------
\lstset{
  language=Python,
  backgroundcolor=\color{codebg},
  frame=single,
  rulecolor=\color{codeframe},
  basicstyle=\ttfamily\small,
  keywordstyle=\color{keywordcol}\bfseries,
  commentstyle=\color{commentcol}\itshape,
  stringstyle=\color{stringcol},
  numberstyle=\tiny\color{numbercol},
  showstringspaces=false,
  breaklines=true,
  breakatwhitespace=true,
  tabsize=4,
  aboveskip=0.8em,
  belowskip=0.8em,
}

% ---- callout boxes ----------------------------------------------------------
\tcbuselibrary{skins, breakable}

\newenvironment{note}{%
  \begin{tcolorbox}[
    colback=notebg, colframe=noteframe,
    fonttitle=\bfseries, title=Note,
    breakable, enhanced
  ]
}{%
  \end{tcolorbox}
}

\newenvironment{intuition}{%
  \begin{tcolorbox}[
    colback=warnbg, colframe=warnframe,
    fonttitle=\bfseries, title=Intuition,
    breakable, enhanced
  ]
}{%
  \end{tcolorbox}
}

% ---- section styling --------------------------------------------------------
\titleformat{\section}{\large\bfseries}{}{0em}{}[\titlerule]
\titleformat{\subsection}{\normalsize\bfseries}{}{0em}{}

% ---- page style -------------------------------------------------------------
\pagestyle{fancy}
\fancyhf{}
\lhead{\textit{ML Library Primer -- NeonTreeEvaluation Working Group}}
\rhead{\thepage}
\renewcommand{\headrulewidth}{0.4pt}

% ---- hyperref ---------------------------------------------------------------
\hypersetup{
  colorlinks=true,
  linkcolor=keywordcol,
  urlcolor=keywordcol,
}

% =============================================================================
\begin{document}

\begin{center}
  {\LARGE\bfseries ML Library Primer}\\[0.4em]
  {\large NeonTreeEvaluation Working Group}\\[0.2em]
  {\normalsize\color{gray} Keep this open while reading the code}
\end{center}

\vspace{1em}
\tableofcontents
\vspace{1em}

% =============================================================================
\section{The Big Picture}

Our task is to automatically detect individual tree crowns in aerial imagery
collected by the National Ecological Observatory Network (NEON). The dataset
includes:

\begin{itemize}[nosep]
  \item \textbf{RGB imagery} -- standard colour photos taken from the air
        (think Google Maps, but scientific quality).
  \item \textbf{LiDAR-derived canopy height} -- a laser rangefinder sweep that
        gives us a height value for every point on the ground, letting us
        distinguish tall trees from shrubs and grass even when everything looks
        green in the photo.
  \item \textbf{Hyperspectral imagery} -- hundreds of light wavelengths beyond
        what the human eye can see, which can distinguish tree species by how
        their leaves reflect light.
  \item \textbf{Labels} -- for each plot, expert annotators drew bounding boxes
        around every visible tree crown. These are our training targets.
\end{itemize}

A machine learning model will look at a patch of this imagery and output a
list of boxes (and possibly pixel-level masks) indicating where it thinks tree
crowns are. To build, train, and run such a model in Python we need three
libraries: \textbf{NumPy}, \textbf{PyTorch}, and \textbf{Torchvision}. This
document explains what each one does and why we reach for it.

\begin{note}
  These libraries do not know anything about trees. They are general-purpose
  mathematical and engineering tools that happen to be the lingua franca of
  computer vision research. You do not need to understand the mathematics
  deeply to work with them -- but you do need to understand their data
  representations, because every bug you will encounter ultimately comes down
  to a misunderstanding about how data is shaped or stored.
\end{note}

% =============================================================================
\section{NumPy: Images Are Grids of Numbers}

\subsection{Why a separate library at all?}

Python lists can store numbers, so why not just use those? The problem is
speed. A single RGB image tile from the NEON dataset might be $400 \times 400$
pixels, meaning 480\,000 numbers. A Python list of 480\,000 numbers is slow
because Python checks the type of each number individually every time you do
anything with it. NumPy avoids this by storing all the numbers in a contiguous
block of memory and running operations on the whole block at once in compiled C
code. The difference is typically 100--1000x faster.

\subsection{The array: NumPy's one big idea}

NumPy has essentially one data structure: the \emph{array} (called
\lstinline{ndarray} internally). Everything else is operations on arrays. An
array is a grid of numbers, all of the same type, with a fixed shape.

\begin{lstlisting}
import numpy as np

# A 1-D array -- just a list of numbers
heights = np.array([15.2, 8.7, 22.1, 5.0])

# A 2-D array -- a grid (like a spreadsheet)
patch = np.zeros((400, 400))  # 400 rows, 400 columns, all zeros

# A 3-D array -- what a single RGB image looks like
rgb_image = np.zeros((400, 400, 3))  # height x width x channels
\end{lstlisting}

The shape of an array is the most important thing to keep in your head at all
times. \lstinline{rgb_image.shape} returns \lstinline{(400, 400, 3)}, which
means 400 rows, 400 columns, and 3 values per pixel (red, green, blue, each
ranging from 0 to 255).

\subsection{How images become arrays}

Every image you have ever seen on a screen is a grid of pixels. Each pixel is
just a few numbers. An RGB pixel is three numbers: how much red, how much
green, how much blue. A single NEON RGB tile therefore becomes a
\lstinline{(height, width, 3)} array of integers.

The LiDAR canopy height raster is a \lstinline{(height, width, 1)} array --
one number per pixel, representing height above ground in metres. You will
sometimes see this stored as \lstinline{(height, width)} (no trailing
dimension) -- these two shapes refer to the same physical data; the only
difference is notational.

The hyperspectral data has 369 wavelength bands, making each tile a
\lstinline{(height, width, 369)} array. That is a lot of channels, which is
part of why ML approaches are appealing -- no human can manually reason about
369 simultaneous measurements per pixel.

\begin{intuition}
  Think of a 3-D array as a stack of 2-D grids, one per channel. For an RGB
  image: layer 0 is the ``red intensity'' map, layer 1 is ``green intensity'',
  layer 2 is ``blue intensity''. For hyperspectral data: layer 0 is how
  brightly the scene reflects 380\,nm ultraviolet light, layer 1 is 382\,nm,
  and so on up to near-infrared.
\end{intuition}

\subsection{Slicing: picking out parts of an array}

You will frequently need to crop a region out of an image, or isolate one
channel. NumPy uses the same slice syntax as Python lists, extended to multiple
dimensions:

\begin{lstlisting}
# Load a NEON tile (rasterio is the library that reads GeoTIFF files;
# it hands the data back to you as a numpy array)
import rasterio
with rasterio.open("data/RGB/SJER_006.tif") as src:
    tile = src.read()  # shape: (3, height, width) -- note: channels first!

# Crop to a 256x256 patch starting at pixel (100, 200)
patch = tile[:, 100:356, 200:456]
# tile[:, ...]  means "all channels"
# 100:356       means rows 100 up to (but not including) 356
# 200:456       means columns 200 up to (but not including) 456

# Pull out just the red channel
red_channel = tile[0, :, :]   # shape: (height, width)

# Pull out just the green channel
green_channel = tile[1, :, :] # shape: (height, width)
\end{lstlisting}

\begin{note}
  Notice that \lstinline{rasterio} returns images as
  \lstinline{(channels, height, width)} -- channels first -- while the
  convention you will see elsewhere (e.g. when building arrays by hand) is
  often \lstinline{(height, width, channels)} -- channels last. This
  inconsistency is a common source of bugs. Always check
  \lstinline{array.shape} when loading data from a new source.
\end{note}

\subsection{Vectorised operations: no for loops over pixels}

The most important habit to develop with NumPy is to avoid writing Python
\lstinline{for} loops over array elements. If you find yourself writing
\lstinline{for i in range(height): for j in range(width): ...}, there is
almost certainly a NumPy function that does it in a single line, thousands of
times faster.

\begin{lstlisting}
# Bad: slow Python loop
for i in range(400):
    for j in range(400):
        tile[0, i, j] = tile[0, i, j] / 255.0

# Good: vectorised -- the division applies to every element at once
tile = tile / 255.0

# Normalise each channel to zero mean and unit variance
mean = tile.mean(axis=(1, 2), keepdims=True)  # shape: (3, 1, 1)
std  = tile.std(axis=(1, 2), keepdims=True)
tile_normalised = (tile - mean) / std
\end{lstlisting}

The \lstinline{axis} parameter tells NumPy which dimensions to collapse. In
the example above, we want one mean per channel, so we average over axes 1 and
2 (height and width) while keeping axis 0 (channels) intact.

% =============================================================================
\section{PyTorch: NumPy That Can Learn}

\subsection{Why not just use NumPy?}

NumPy is excellent for manipulating data, but it cannot train a neural network
on its own. There are two things it is missing:

\begin{enumerate}[nosep]
  \item \textbf{GPU support.} Training a modern object detector on a dataset
        the size of NeonTreeEvaluation on a CPU would take weeks. A GPU
        performs the same arithmetic but has thousands of cores running in
        parallel, reducing this to hours. NumPy has no GPU support.
  \item \textbf{Automatic differentiation.} Training a neural network means
        repeatedly computing ``how should each parameter change to reduce the
        error?'' This requires calculating gradients -- derivatives of the
        error with respect to every parameter. Doing this by hand for a model
        with millions of parameters is impossible. PyTorch does it
        automatically.
\end{enumerate}

PyTorch solves both problems. Its fundamental data structure is the
\emph{tensor}, which is essentially a NumPy array that can live on a GPU and
can track gradients.

\subsection{Tensors}

A tensor behaves almost identically to a NumPy array:

\begin{lstlisting}
import torch
import numpy as np

# Create a tensor directly
t = torch.zeros((3, 400, 400))   # same syntax as np.zeros

# Convert from numpy (shares memory -- no data is copied)
arr = np.random.rand(3, 400, 400).astype(np.float32)
t = torch.from_numpy(arr)

# Convert back to numpy
arr_back = t.numpy()

# Arithmetic works the same way
t_normalised = (t - t.mean()) / t.std()
print(t_normalised.shape)   # torch.Size([3, 400, 400])
\end{lstlisting}

The main practical difference you will notice is the method names. NumPy uses
\lstinline{np.zeros}, PyTorch uses \lstinline{torch.zeros}. NumPy uses
\lstinline{arr.shape}, PyTorch also uses \lstinline{t.shape} but wraps the
result in a \lstinline{torch.Size} object (which behaves like a tuple).

\subsection{Moving data to the GPU}

When a GPU is available, you move tensors to it explicitly:

\begin{lstlisting}
device = torch.device("cuda" if torch.cuda.is_available() else "cpu")
print(device)  # "cuda" on a GPU machine, "cpu" otherwise

t = t.to(device)  # move tensor to GPU (or keep on CPU if no GPU)
\end{lstlisting}

Once on the GPU, all operations on that tensor happen on the GPU
automatically. The catch is that you cannot mix GPU and CPU tensors in the
same operation -- both operands must be on the same device. This is probably
the most common runtime error you will encounter:
\lstinline{RuntimeError: Expected all tensors to be on the same device}.

\subsection{Automatic differentiation}

This is the feature that makes PyTorch indispensable for training. Every
tensor can optionally record a history of the operations applied to it so that
PyTorch can later compute gradients automatically. You will not need to
implement this yourself -- the training loop does it for you -- but
understanding what is happening underneath prevents confusion.

\begin{lstlisting}
# Model parameters need gradient tracking
weight = torch.tensor([0.5], requires_grad=True)

# A trivial "model": prediction = weight * input
x = torch.tensor([2.0])
prediction = weight * x   # = 1.0

# Suppose the true answer is 3.0; compute squared error
loss = (prediction - 3.0) ** 2   # = 4.0

# Ask PyTorch to compute d(loss)/d(weight)
loss.backward()
print(weight.grad)   # tensor([-8.]) -- loss decreases if we increase weight
\end{lstlisting}

\lstinline{loss.backward()} is the single call that triggers PyTorch to
differentiate the entire computation graph. In a real model with millions of
parameters, the same call computes all the gradients simultaneously.

\subsection{The training loop}

In practice you will see the same four-step pattern repeated for every batch
of training data:

\begin{lstlisting}
optimizer = torch.optim.Adam(model.parameters(), lr=1e-4)

for images, targets in dataloader:
    images = images.to(device)
    targets = [{k: v.to(device) for k, v in t.items()} for t in targets]

    # 1. Forward pass: run the model, get predictions and loss
    loss_dict = model(images, targets)
    loss = sum(loss_dict.values())

    # 2. Zero gradients from the previous batch (must do this every iteration)
    optimizer.zero_grad()

    # 3. Backward pass: compute gradients
    loss.backward()

    # 4. Update parameters in the direction that reduces loss
    optimizer.step()
\end{lstlisting}

The optimizer (\lstinline{torch.optim.Adam} here) implements the rule for
updating parameters given their gradients. Adam is a standard choice that
works well across a wide range of problems without much tuning.

\subsection{Datasets and DataLoaders}

PyTorch has a pair of abstractions for feeding data to a model in a structured
way.

A \textbf{Dataset} knows how to return one example (one image and its
annotations) given an index. You implement one by subclassing
\lstinline{torch.utils.data.Dataset} and defining two methods:

\begin{lstlisting}
class NeonDataset(torch.utils.data.Dataset):
    def __init__(self, image_paths, annotation_paths, transform=None):
        self.image_paths = image_paths
        self.annotation_paths = annotation_paths
        self.transform = transform

    def __len__(self):
        # How many examples are in this dataset?
        return len(self.image_paths)

    def __getitem__(self, idx):
        # Load and return the idx-th example
        image = load_image(self.image_paths[idx])    # returns a tensor
        boxes = load_boxes(self.annotation_paths[idx]) # returns a tensor of [x1,y1,x2,y2]

        if self.transform:
            image = self.transform(image)

        target = {"boxes": boxes, "labels": torch.ones(len(boxes), dtype=torch.long)}
        return image, target
\end{lstlisting}

A \textbf{DataLoader} wraps a Dataset and handles batching, shuffling, and
parallel loading:

\begin{lstlisting}
dataset = NeonDataset(image_paths, annotation_paths, transform=preprocess)
dataloader = torch.utils.data.DataLoader(
    dataset,
    batch_size=4,        # process 4 images at once
    shuffle=True,        # randomise order each epoch
    num_workers=4,       # load data in parallel on 4 CPU workers
)

# The training loop iterates over the DataLoader
for images, targets in dataloader:
    ...  # images is now a batch of 4 tensors stacked together
\end{lstlisting}

\begin{note}
  When the DataLoader collects individual examples into a batch, the shape of
  each image tensor changes from \lstinline{(C, H, W)} to
  \lstinline{(B, C, H, W)}, where \lstinline{B} is the batch size. This
  leading batch dimension is ubiquitous in PyTorch code; if you see a tensor
  with an unexpected extra leading dimension, it is usually the batch.
\end{note}

% =============================================================================
\section{Torchvision: The Image-Specific Toolkit}

Torchvision is an official PyTorch companion library that provides three
things: a collection of image transforms, pre-trained model architectures, and
utilities for common vision datasets. We care primarily about the first two.

\subsection{Transforms}

Raw imagery cannot be fed directly into most models. The images need to be
pre-processed to a consistent size, normalised so that pixel values have
similar magnitudes, and sometimes augmented (randomly flipped, cropped, etc.)
to help the model generalise. Torchvision provides a composable set of
transforms for this purpose.

\begin{lstlisting}
import torchvision.transforms.v2 as T

# A preprocessing pipeline applied to every image before it enters the model
preprocess = T.Compose([
    T.ToImage(),                         # convert numpy array to torch tensor
    T.ToDtype(torch.float32, scale=True),# convert uint8 [0,255] to float [0,1]
    T.Resize((512, 512)),                # resize to a fixed resolution
    T.Normalize(
        mean=[0.485, 0.456, 0.406],      # ImageNet channel means
        std= [0.229, 0.224, 0.225],      # ImageNet channel stds
    ),
])

# A data augmentation pipeline used only during training
augment = T.Compose([
    T.RandomHorizontalFlip(p=0.5),
    T.RandomVerticalFlip(p=0.5),         # aerial imagery has no canonical top
    T.ColorJitter(brightness=0.2, contrast=0.2),
])
\end{lstlisting}

\begin{intuition}
  The \lstinline{Normalize} step uses the channel means and standard
  deviations computed from the ImageNet dataset (a large collection of
  everyday photographs). We reuse these even for aerial imagery because most
  of our model's weights will have been pre-trained on ImageNet data (see
  below), and normalising with the same statistics the pre-training used keeps
  the input distribution consistent.
\end{intuition}

The \lstinline{T.Compose} call chains transforms together so they run in
sequence. The resulting object is callable -- you apply it to an image like a
function: \lstinline{processed = preprocess(raw_image)}.

For object detection tasks, transforms need to be applied to both the image
and the bounding boxes simultaneously (a random crop that removes part of the
image must also update the box coordinates). The \lstinline{v2} version of
Torchvision transforms handles this automatically when you pass the target
dictionary alongside the image.

\subsection{Pre-trained models and transfer learning}

Training a neural network from random initial weights on a dataset as small as
NeonTreeEvaluation (around 6000 annotated trees) would not work well. There is
not enough data for the model to learn low-level visual features like
``edge'', ``texture'', and ``blob of similar colour'' from scratch. Instead,
we use \emph{transfer learning}: start from a model that has already learned
these features on a large dataset, then fine-tune it on our data.

Torchvision provides this out of the box:

\begin{lstlisting}
from torchvision.models.detection import fasterrcnn_resnet50_fpn, FasterRCNN_ResNet50_FPN_Weights

# Load Faster R-CNN pre-trained on COCO (a large object detection dataset)
weights = FasterRCNN_ResNet50_FPN_Weights.DEFAULT
model = fasterrcnn_resnet50_fpn(weights=weights)

# Replace the final classification head to detect only one class (tree crown)
# The rest of the network keeps its pre-trained weights
in_features = model.roi_heads.box_predictor.cls_score.in_features
from torchvision.models.detection.faster_rcnn import FastRCNNPredictor
model.roi_heads.box_predictor = FastRCNNPredictor(in_features, num_classes=2)
# num_classes=2 means: background (class 0) + tree crown (class 1)

model = model.to(device)
\end{lstlisting}

The pre-trained backbone (ResNet-50) has learned to recognise textures and
shapes across millions of photographs. We discard only the final layer that
maps features to COCO category labels, replace it with a fresh layer for our
single class, and then train the whole thing on our aerial imagery. The
backbone will update its weights slowly to adapt to the overhead perspective
and tree-specific appearance while retaining the general visual understanding
it was initialised with.

\subsection{Inference: running the model on new data}

During inference (applying the model to images it has not seen before) there
is no need to track gradients, which saves memory and time. PyTorch provides a
context manager for this:

\begin{lstlisting}
model.eval()   # switch model to evaluation mode (disables dropout etc.)

with torch.no_grad():
    for images, _ in test_dataloader:
        images = [img.to(device) for img in images]
        predictions = model(images)
        # predictions is a list of dicts, one per image:
        # {"boxes": tensor of shape (N, 4),
        #  "labels": tensor of shape (N,),
        #  "scores": tensor of shape (N,)}

        for pred in predictions:
            keep = pred["scores"] > 0.5          # filter low-confidence detections
            boxes = pred["boxes"][keep].cpu().numpy()
            # boxes is now a numpy array of [x1, y1, x2, y2] coordinates
\end{lstlisting}

Note the \lstinline{.cpu().numpy()} chain: we move the tensor back to CPU
memory (from the GPU) and then convert to a NumPy array so we can work with it
using standard tools.

% =============================================================================
\section{The Full Pipeline in One View}

To close, here is a sketch of how the pieces connect, from raw data on disk to
a trained model producing predictions:

\begin{enumerate}
  \item \textbf{Load} -- \lstinline{rasterio} reads a GeoTIFF tile and returns
        a NumPy array of shape \lstinline{(channels, H, W)}.
  \item \textbf{Package} -- a custom \lstinline{NeonDataset} wraps the file
        paths and knows how to load one image and its bounding-box annotations
        on demand.
  \item \textbf{Preprocess} -- Torchvision transforms convert the array to a
        float tensor, resize it, and normalise the pixel values.
  \item \textbf{Batch} -- a \lstinline{DataLoader} collects 4 (or 8, or 16)
        processed images into a batch of shape \lstinline{(B, C, H, W)}.
  \item \textbf{Forward pass} -- the batch goes into the model (Faster R-CNN
        or similar); PyTorch runs the computation on the GPU.
  \item \textbf{Loss} -- the model compares its predicted boxes against the
        ground-truth annotations and returns a scalar loss value.
  \item \textbf{Backward pass} -- \lstinline{loss.backward()} computes
        gradients automatically.
  \item \textbf{Update} -- the optimizer nudges every parameter slightly in the
        direction that reduces the loss.
  \item \textbf{Repeat} -- steps 4--8 repeat for every batch, for multiple
        passes through the full dataset (called \emph{epochs}).
\end{enumerate}

The forestry team is most likely to be involved in steps 1--2 (loading and
inspecting data, verifying that annotations look correct, splitting the dataset
into train and test sets) and step 9 (deciding when the model's predictions
are good enough for ecological conclusions). Steps 5--8 largely run on their
own once the scaffold is in place.

\begin{note}
  A recurring theme in this document is that you should always know the
  \emph{shape} of every array or tensor you are working with. When code
  breaks, the first thing to do is print the shapes of the tensors involved.
  Most errors in ML code come down to a shape mismatch: two tensors that were
  expected to be compatible turned out to have different sizes along some
  dimension.
\end{note}

\end{document}
